\chapter{Нотация}
\label{ch:notation}

Сводка обозначений, используемых по тексту.

\begin{longtable}{p{0.20\textwidth} p{0.74\textwidth}}
\toprule
\textbf{Символ} & \textbf{Смысл} \\
\midrule
\endhead
$G=(V,E)$ & граф памяти: узлы $V$, направленные типизированные знаковые рёбра $E$ \\
$v,\,u$ & узлы (воспоминания, факты, сущности, концепты) \\
$\Neigh(A)$ & окрестность узла $A$ (соседи, фронт распространения) \\
$\mathbf{x}_v\in\mathbb{R}^d$ & векторное представление содержания узла (эмбеддинг) \\
$c_v=\conf(v)\in[-1,1]$ & знаковая уверенность: $+$да, $-$нет, $0$ незнание \\
$a_v$ & транзиентная активация узла при обходе \\
$\text{importance}(v)$ & важность узла $\in [0,1]$: salience при кодировании + центральность; гейтит прунинг \\
$r_{uv}$ & тип отношения на ребре $u\to v$; знак $\pm$ \\
$w_{uv}$ & вес (эвиденциальная сила) ребра \\
$m_{u\to v}$ & сообщение, передаваемое по ребру: $m_{u\to v}=\phi_{r_{uv}}(h_u,h_v,e_{uv})$ \\
$\tau$ & постоянная времени слоя (интегратора) \\
$B_v$ & базовая активация ACT-R узла $v$ \\
$\pi$ & стационарное распределение случайного блуждания \\
$P$ & матрица переходов блуждающего агента; $P,Q$ --- операторы проекции/поднятия при огрублении \\
$\alpha$ & доля телепортов (рестартов) в блуждании \\
$L,\,L^{+}$ & лапласиан графа и его псевдообратный \\
$\Reff(i,j)$ & эффективное сопротивление между узлами $i,j$ \\
$\sigma$ & ветвящийся коэффициент лавины пропагации \\
$\theta_{\text{emergent}},\theta_{\text{temp}}$ & верхний и нижний пороги распространения (гистерезис) \\
$\Delta\conf(B\mid A)$ & локальное изменение уверенности соседа $B$ под влиянием факта $A$ \\
$\mood_t\in\mathbb{R}^k$ & вектор настроения в момент $t$ (аффективное пространство) \\
$\varphi(\cdot)\in\mathbb{R}^d$ & скорер по осям характера \\
$\mathbf{t}$ & целевые координаты Рин в пространстве характера \\
$\mathrm{Beta}(\alpha,\beta)$ & апостериорное распределение успешности реакции \\
$\eta$ & темп обновления (learning rate) локального правила \\
\bottomrule
\end{longtable}
